\section{Desenvolvimento}

\subsection{Topico 1: primeiro aspecto do tema}

\topicclaim{Explique aqui a ideia central deste primeiro topico.}

Desenvolva o primeiro aspecto que esta sob o mesmo guarda-chuva tematico da
tese. Apresente evidencias, exemplos, conceitos ou autores que ajudem o leitor
a entender por que este ponto e relevante para o argumento principal.

Ao final do paragrafo ou da subsecao, conecte explicitamente este topico a tese
do essay. Isso evita que o desenvolvimento vire apenas uma lista de assuntos.

Quando o essay precisar apresentar um procedimento, use um algoritmo numerado e
referenciavel, como no Algoritmo~\ref{alg:busca-completa}. O exemplo tambem
mostra uma forma compacta de registrar variaveis, condicoes, retorno,
pos-condicao e complexidade.

\begin{algorithm}[H]
\caption{Busca linear com sentinela de controle}
\label{alg:busca-completa}
\Entrada{Vetor $A[1 \ldots n]$, tamanho $n \in \mathbb{N}$ e chave $x$}
\Saida{Indice $p$ tal que $A[p] = x$, ou $-1$ caso $x \notin A$}

$p \gets -1$\tcp*{valor padrao quando a chave nao e encontrada}
$i \gets 1$\tcp*{variavel de iteracao}
$comparacoes \gets 0$\tcp*{contador auxiliar}

\While{$i \leq n$ \textbf{e} $p = -1$}{
  $comparacoes \gets comparacoes + 1$\;

  \eIf{$A[i] = x$}{
    $p \gets i$\tcp*{guarda a primeira ocorrencia}
  }{
    $i \gets i + 1$\;
  }
}

\tcp{Pos-condicao: $p = -1$ ou $A[p] = x$}
\tcp{Complexidade: $T(n) \in O(n)$ no pior caso e $T(n) \in \Omega(1)$ no melhor caso}
\tcp{Espaco adicional: $S(n) \in O(1)$}
\Retorne $p$\;
\end{algorithm}

No texto, as variaveis matematicas aparecem em modo matematico, como $A$, $n$,
$x$, $i$ e $p$. A atribuicao e escrita com $\gets$; intervalos podem ser
representados por $1 \ldots n$; pertencimento e ausencia por $\in$ e $\notin$;
e a analise assintotica usa $O(n)$, $\Omega(1)$ ou $\Theta(n)$, conforme o caso.
